\documentclass[10pt,a4paper]{article}
\usepackage[T1]{fontenc}
\usepackage{lmodern}
\usepackage[margin=23mm]{geometry}
\usepackage{amsmath,amssymb,booktabs,microtype,tabularx}
\usepackage[hidelinks,hypertexnames=false]{hyperref}
\pdfinfoomitdate=1
\pdftrailerid{}
\pdfsuppressptexinfo=15
\setlength{\parindent}{0pt}
\setlength{\parskip}{5pt}
\begin{document}
\begin{center}
{\Large A stellar spectrum as prescribed hydrogen irradiation}\\[4pt]
\texttt{exojax-escape}: model definition and limits, 23 November 2026
\end{center}

\section{From surface flux to photon occupation}
The new input replaces the blackbody spectrum that excites and ionizes the
planetary gas. It does not solve radiation transport through that gas.
The public PHOENIX proxy has $T_{\rm eff}=10\,000$ K, $\log g=4.0$ and
solar composition. It is a synthetic LTE stellar atmosphere, not a measured
far-UV spectrum or the recovered exact stellar input of Zhang et al.

The library supplies outward-hemisphere surface flux
$F_\lambda^{\star}$ in erg\,s$^{-1}$\,cm$^{-2}$\,cm$^{-1}$ and vacuum
wavelength in \AA. Converting wavelength to cm first gives
\begin{equation}
 \nu=\frac{c}{\lambda},\qquad
 F_\nu^{\star}=F_\lambda^{\star}\frac{\lambda^2}{c},\qquad
 J_\nu=W\frac{F_\nu^{\star}}{\pi},\qquad
 \eta_\nu=\frac{c^2J_\nu}{2h\nu^3}
         =\frac{Wc^2F_\nu^{\star}}{2\pi h\nu^3}.
 \tag{SED-1}\label{eq:conversion}
\end{equation}
$F_\nu^{\star}$ is per Hz; $J_\nu$ is mean intensity per steradian;
$\eta_\nu$ is dimensionless. The factor $\pi$ assumes a uniformly bright
stellar disk, $I_\nu^{\star}=F_\nu^{\star}/\pi$. With disk solid angle
$\Omega_\star$, $W=\Omega_\star/(4\pi)$; the example retains $W=0.025$.
There is no second dilution or luminosity normalization. Substituting
$F_\nu^{\star}=\pi B_\nu(T)$ recovers
$\eta_\nu=W/[\exp(h\nu/k_{\rm B}T)-1]$.

\path{src/escape/stellar_radiation.py:surface_flux_occupation} linearly
interpolates $F_\nu^{\star}$ on a strictly increasing frequency grid and
evaluates Eq.~\eqref{eq:conversion}. Invalid values and frequencies outside
coverage yield NaN; they are not silently endpoint-clamped. Gradients in
surface flux and dilution pass through this conversion.

\section{How the spectrum enters the atom}
Both arrays are required. They sample the atom's transition frequencies and
continuum quadrature. Pass them to the \path{population_profiles} or
\texttt{spectrum} method of \texttt{HydrogenNLTEForwardModel}.
They already include dilution. Otherwise the previous
diluted-blackbody path remains the default. Each layer $r$ solves
\begin{equation}
 (\{n_i(r)\},n_e(r),n_p(r))
 =\mathcal S\!\left[n_{\rm H}(r),T(r);
       \{\eta^{\rm line}_{ul}\},\{\eta^{\rm cont}_q\}\right],
 \qquad \eta^{\rm line}_{ul},\eta^{\rm cont}_q\ \hbox{independent of }r.
 \tag{SED-2}\label{eq:shared}
\end{equation}
$\mathcal S$ denotes the existing charge-conserving local statistical
equilibrium solve. The density, kinetic temperature and Parker velocity
remain prescribed. Ionization is recomputed; it does not feed back into the
wind pressure or an energy equation. Existing branch opacities/emissivities
then generate the observer spectrum. The incident stellar spectrum does not
replace the observer's stellar reference or observation-processing model.

The default example samples each stellar line at the atomic center.
For optional prescribed velocity/smoothing sensitivities it instead uses
\begin{equation}
 \eta^{\rm line}_{ul}=\sum_a w_a\,
 \eta_{\nu_{ul}[1+(v_a+\Delta v)/c]},\qquad \sum_a w_a=1.
 \tag{SED-3}\label{eq:averaging}
\end{equation}
Uniform-disk rotation and a fixed Gaussian kernel only test sensitivity to
spectral sampling. They are not angle- and layer-dependent line-profile
averages. The stellar rotation projected toward the planet need not equal
the observer's $v\sin i$. Continuum quadrature samples the unbroadened input.
Positive \texttt{shift\_kms} in Eq.~\eqref{eq:averaging} samples higher
stellar frequency (bluer wavelength). It is a sampling offset, not an
observer radial velocity or a solved star--planet velocity geometry.
The stellar-absorption plot instead uses $\nu=\nu_0(1-v/c)$, so its positive
horizontal velocity denotes redder wavelength.

\newpage
\section{Coverage and the controlled comparison}
PHOENIX covers 500--54\,999.75 \AA\ in this file. The example explicitly
completes both tails with a $10\,000$ K Planck shape matched to the nearest
surface-flux endpoint. Zero and unscaled-blackbody UV tails, and an unscaled
blackbody IR tail, are separate sensitivities. A matched tail is an assumed
boundary continuation, not measured stellar EUV emission. No extra surface
flux normalization is applied.

\path{examples/stellar_halpha.py} compares the full input replacement and
separate line/continuum replacements at two fixed previously specified
atmospheres. It also records line shifts, smoothing, projected rotation and
tail changes. Populations and charge are recomputed for every field. A
constant offset is profiled for the displayed observation comparison:
\begin{equation}
 b=\frac{\sum_k(d_k-f_k)/\sigma_k^2}{\sum_k1/\sigma_k^2},\qquad
 R=\left[\frac{1}{N}\sum_k
       \left(\frac{f_k+b-d_k}{\sigma_k}\right)^2\right]^{1/2}.
 \tag{SED-4}\label{eq:comparison}
\end{equation}
$R$ is a descriptive normalized residual RMS, not a validated likelihood or
a posterior result. No physical parameters are refitted in this experiment.
Keep numerical refinement errors separate from physical model error.

\section{What absorbs Ly\texorpdfstring{$\boldsymbol\beta$}{beta}?}
The photospheric absorption already present in the stellar input is distinct
from absorption by the planet. In hydrogen, Ly$\beta$ couples $1s$ to $3p$.
Ground-state neutral hydrogen can therefore intercept these photons in
either atmosphere. Excited $3p$ can decay back to $1s$ or to $2s$; repeated
resonance scattering and conversion into other photons must be distinguished
from simple photon destruction. Reducing one imposed occupation by hand
does not model this transport.

The missing planetary coupling has the schematic form
\begin{equation}
 \frac{\mathrm d I_\nu}{\mathrm ds}=-\alpha_\nu I_\nu+j_\nu,
 \qquad J_\nu=\frac{1}{4\pi}\int I_\nu\,\mathrm d\Omega,
 \qquad (\alpha_\nu,j_\nu)=(\alpha_\nu,j_\nu)[\{n_i\},T].
 \tag{SED-5}\label{eq:missing}
\end{equation}
Equation~\eqref{eq:missing} must feed back into Eq.~\eqref{eq:shared}; this
experiment does not implement that loop. Frequency redistribution, moving
gas, diffuse emission, the incident UV uncertainty and thermal/ionization
feedback remain physical questions. A realistic stellar photospheric
absorption trough alone does not determine the radiation inside the planet.

\section{What MCEG records}
\href{../mceg/contracts/escape.stellar_radiation.yaml}{The new contract}
binds Eqs.~\eqref{eq:conversion}--\eqref{eq:shared} to their current code,
separates finite numerical tests from atmospheric accuracy, and retains
physical adequacy as \textbf{unknown}. A direct FollowUp requires a matched
radiation-plus-population comparison and a declared physical error budget.
The limited-use Decision permits exploratory sensitivity work only.

Earlier provider contracts retain their historical source hashes and tests;
those tests are not promoted to evidence for the changed implementation.
MCEG remains an external verification tool, with no runtime dependency from
\texttt{escape}. Raw or derived PHOENIX tables are not bundled: the official
library requests citation, but an explicit redistribution license was not
found in the inspected pages.

\small
\textbf{References.}
Husser et al. (2013), A\&A 553, A6,
\href{https://doi.org/10.1051/0004-6361/201219058}{doi:10.1051/0004-6361/201219058};
\href{https://phoenix.astro.physik.uni-goettingen.de/spectra/}{official PHOENIX files and conventions}.
Zhang et al. (2026), \href{https://arxiv.org/abs/2606.23899v1}{arXiv:2606.23899v1},
Section IV.1 (the exact input configuration is unavailable here).
\end{document}
